% Generated draft: requires mathematical and source review.
\documentclass[11pt]{article}
\usepackage[margin=25mm]{geometry}
\usepackage{fontspec}
\setmainfont{Latin Modern Roman}
\usepackage{amsmath,amssymb,mathtools,unicode-math}
\setmathfont{Latin Modern Math}
\usepackage{enumitem,needspace,xurl,hyperref}
\setlength{\parindent}{0pt}
\setlength{\parskip}{5pt}
\setlength{\emergencystretch}{3em}
\newcommand{\N}{\mathbb{N}}
\newcommand{\Z}{\mathbb{Z}}
\newcommand{\Q}{\mathbb{Q}}
\newcommand{\R}{\mathbb{R}}
\newcommand{\C}{\mathbb{C}}
\newcommand{\F}{\mathbb{F}}
\newcommand{\A}{\mathbb{A}}
\newcommand{\PP}{\mathbb{P}}
\newcommand{\E}{\mathbb{E}}
\newcommand{\eps}{\varepsilon}
\newcommand{\dd}{\,\mathrm{d}}
\begin{document}
\section*{UnsolvedMath ID: 30003695}
\textbf{Source problem code:} catalogue entry
\subsection{Definitions and mathematical statement}
Let \(\mathrm{Ab}\) denote the category of abelian groups, and let \(k\in\Z_{\geq 0}\). The supplied record names a category \(S(\Z)\), a family of groups \(\operatorname{IA}_{\bullet}\), and a homology functor
\[
\mathcal H_k:=H_k(\operatorname{IA}_{\bullet})\colon S(\Z)\longrightarrow\mathrm{Ab}.
\]
Here \(H_k\) denotes group homology in degree \(k\), but the coefficient module and its action are not specified in the record.

\textbf{Conjecture attributed to Djament in the supplied curation.} For every nonnegative integer \(k\), the functor \(\mathcal H_k\) is weakly polynomial of degree at most \(3k\):
\[
\forall k\in\Z_{\geq 0},\qquad
\deg_{\mathrm{weak}}\!\left(H_k(\operatorname{IA}_{\bullet})\right)\leq 3k.
\]
The notation \(\deg_{\mathrm{weak}}(F)\leq d\) here abbreviates the source's assertion that \(F\) is weakly polynomial of degree at most \(d\).

\textbf{Unresolved definitions.} The supplied record does not define the objects or morphisms of \(S(\Z)\), the groups \(\operatorname{IA}_{\bullet}\), the maps making their homology a functor on this category, or the intended definition of weak polynomiality. A customary possible interpretation of \(\operatorname{IA}_n\) is the kernel of the action of the automorphism group of a free group of rank \(n\) on its abelianization, but that interpretation is not specified by the supplied text and is not imposed here. Likewise, integral homology with trivial coefficients is not silently assumed. The displayed conjecture therefore preserves the source formulation, but requires these definitions and conventions to be confirmed before it becomes a self-contained mathematical statement.
\subsection{Short English statement}
The supplied conjecture, attributed to Djament, asks whether the weak polynomial degree of the IA-group homology functor in every degree \(k\geq 0\) is at most \(3k\). The category, IA groups, homology coefficients, functorial maps, and weak polynomiality convention require confirmation from the cited sources.
\subsection{Sources}
\begin{itemize}
\item Supplied local UnsolvedMath registry record, permanent ID \(30003695\), title \emph{Polynomial Degree of IA-Group Homology Functors}; Oberwolfach Reports Open Problems collection; exact source problem code \texttt{catalogue entry}. The record lists proposed year 2018 and supplies the attribution to Djament through its background and literature assessment; its separate proposer field is empty.
\item Supplied source citation: \emph{Topology of Arrangements and Representation Stability (2019)}, Oberwolfach Reports, pp.\ 43--123, DOI 10.4171/owr/2018/2. \url{https://doi.org/10.4171/owr/2018/2}. This is also the supplied literature-source link labelled ``Oberwolfach report.''
\item Supplied literature source: ``Djament and Vespa, weak polynomial functors.'' \url{https://doi.org/10.1093/imrn/rnx099}.
\item Literature assessment dated exactly \texttt{2026-08-21}: the supplied curation labels the problem open, identifies the extracted assertion as Djament's conjecture, and reports that the searched literature did not establish the asserted weak polynomial degree in all homological degrees. This is supplied dated curation, not an independently verified assessment. The record expressly requests independent verification of time-sensitive status and cited claims.
\item The supplied statement-status label is \texttt{corrected\_verified}; its accompanying explanation says that mathematical content was preserved while notation, typography, or omitted definitions were repaired. That label is source metadata, not independent mathematical verification of this draft.
\end{itemize}
\end{document}
